\chapter{Integration notes}

\section{Integration into a Project}
\ifnum \DocumentGroup=1{
There is no need for standard installation procedure for RTA-SUM DIAG. Once RTA-SUM plugin is installed, RTA-SUM DIAG gets available for being integrated into a project. Create a new RTA-CAR project in ISOLAR-B and select the Model Year (MY) version. Open the RTA Code Generator dialog and enable RTA-SUM DIAG by choosing the MY selected during the project creation. If the option \textit{Not Configured} is chosen instead, the module will not be integrated. \\
Please refer to RTA-SUM plugin Getting Started Guide for more details about how to integrate this module into a project.
}
\else{
There is no need for standard installation procedure for DIAG. The folder itself contains all the necessary
artifacts to integrate into a project. Dropping the complete folder structure into a BSW configuration
project will enable the user to configure DIAG and generate all the necessary artifacts.\\
}
\fi 
Refer to \href {https://rtahotline.etas.com/confluence/display/RH/RTA-SUM} {RTA-SUM Knowledge Base} for more details about how to integrate this module into a project.
\section{Module Dependencies}
\label{sec:dependencies}
\begin{itemize}
	\item Function \code{Diag\_SetProgConditions} set initialReflashFlag to TRUE and write it into NvM through interface \code{DIAG\_SSC\_InitialReflashFlag}.
	\item MEC Initial value is taken from the configured value of parameter “MecInitialValue” in DiagGeneral container.
	\item In order to fulfill the requirement to store the Tester Source Address in the \code{DiagnosticSession}\break\code{Control\_Confirmation}, \code{CommunicationControl\_Confirmation} and \code{Dcm\_SetProgConditions} \\ (Dcm callout) routine requests from DCM, AUTOSAR typed Per Instance Memory is used to store this data in non-volatile memory.
	\item When connecting the Dcm to the port interface \code{DataServices\_DID\_F1A0\_Manufacturers}\break\code{EnableCounter\_ManufacturersEnableCounter\_0\_0} of the DIAG module, the definition of the provided \code{ClientServer} \code{Operations} must be verified.
	\item The DIAG module has a dependency on the Dcm module configuration. \\
	Under \code{DcmGeneral}, \code{DcmStoringEnabled} must be set to \code{True}.	
	This will ensure that the function prototype for the Dcm callout function \code{Dcm\_SetProgConditions} is included in \file{DcmCore\_DslDsd\_Cfg.h}.
	The call to the function \code{Diag\_SetProgConditions} implemented in SUM\_DIAG must be added into the \code{Dcm\_SetProgConditions}.
	\item The SUM\_DIAG\_DcmGetRxPermission callout shall be invoked by DCM with ProtocolType and DcmRxPduId to perform protocol pre-emption enable checks; it returns \code{E\_OK} only when protocol, tester source address, and vehicle power mode are valid, otherwise \code{E\_NOT\_OK}. Ensure the DCM callout registration and tester source address configuration align with DIAG expectations. The ECU integrator shall call SUM\_DIAG\_DcmGetRxPermission callout function inside DcmAppl\_DcmGetRxPermission integration code.

\end{itemize}

\section{Configuring DIAG}
Once the complete package has been added to a BSW project, an additional module DIAG should appear.
From here you can open and begin configuring the module.
Specific description of each configurable parameter is available in the “Description” view of
\ifnum \DocumentGroup=1{
ISOLAR-A/B
}
\else
{ISOLAR-A with BCT add-on}
\fi
when a parameter is selected. \\
Refer to Section \ref{sec:configuration} for more details about RTA-SUM DIAG configuration.

\section{Generating Artifacts}
\label{sec:Artifacts}
Generating the C/H/SWCD files from your configuration is similar to other BSW modules delivered from ETAS.
The Code Generation perspective of the 
\ifnum \DocumentGroup=1{
ISOLAR-A/B
}
\else
{ISOLAR-A with BCT add-on}
\fi
should present the user with a dialog to select specific module within the project to generate.
\ifnum \DocumentGroup=1{}
\else{Either a Partial Build or Full Build can be used, depending on the needs of the user.}
\fi
In case of any configuration errors present in DIAG, the console window will output a description of any errors found during generation.
Further, when the user is configuring the module, the Problems Log will show any active errors in the configuration of DIAG.
A third visual aid is the appearance of red X’s in the configuration view over items which are invalid.
Simply hovering over these X’s will provide the user a description of the problem.
These are static validations. \\

After generating the SUM files, execute the MemMap tool to generate memory mapping file of the module (\file{SUM\_DIAG\_MemMap.h}).

\section{Configuring the RTE}
After the user has configured and generated the artifacts for DIAG, the output SWCD file \linebreak(\file{SUM\_Diag\_Cfg\_SWCD.arxml}) can be used for the RTE configuration.
This is a complete SWC description per the AUTOSAR standard of all the users configured ports and interfaces required for the project specific configuration.
This SWCD is a standard ASW component and not a module description.
Therefore, there are no special considerations to be taken into account when configuring the RTE.
It is recommended that if specific changes are required to the SWCD, the user not manually edit this file, rather update and regenerate the configuration with the required changes.

Note: RTE command used for RTE generation

\begin{rtecommandline}
	-nts --samples=memmap --os-define-osenv=RTAOS40 --exclusive-area-optimization=disable --os-output-param=all --strict-unconnected-rport-check=off
\end{rtecommandline}

